%============================================================
% CH1 · P5 · PROBLEM 2 — LONGEST CYCLE IN A GRAPH (LeetCode 2360)
%============================================================
\subsection{Problem 1.5.2: Longest Cycle in a Graph}
\problemlabel{1.5.2}{ssec:ch1-longest-cycle}

\problemcard{1.5.2}{Longest Cycle in a Graph}{LeetCode 2360}{https://leetcode.com/problems/longest-cycle-in-a-graph/}{\exhard\ \lc{Hard}}{Out-degree $\le 1$: walk from every unvisited vertex with global visit times; a walk that returns into itself measures a cycle}

\subsubsection*{1. Problem statement}

A directed graph has $n$ vertices $0..n-1$, and every vertex has \textbf{at most one} outgoing edge: \texttt{edges[i]} is the target of the edge leaving $i$, or $-1$ if there is none. Return the length of the \textbf{longest cycle}, or $-1$ if the graph has no cycle.

\textbf{Constraints:} $2 \le n \le 10^5$, \; $-1 \le \texttt{edges}[i] < n$, \; $\texttt{edges}[i] \ne i$ (no self-loops).

\begin{table}[H]
\centering\small
\begin{tabular}{@{}lcl@{}}
\toprule
\textbf{edges} & \textbf{Output} & \textbf{Explanation} \\
\midrule
$[3,3,4,2,3]$ & 3 & the only cycle is $2 \to 4 \to 3 \to 2$ \\
$[2,-1,3,1]$ & $-1$ & $0 \to 2 \to 3 \to 1 \to$ (nothing): no cycle \\
\bottomrule
\end{tabular}
\end{table}

\begin{diagrambox}
\begin{tikzpicture}
    \begin{scope}
        \node[font=\small\bfseries] at (1.2,1.7) {Example 1: $[3,3,4,2,3]$};
        \foreach \v/\x/\y/\s in {0/-0.6/1/gv, 1/-0.6/-0.6/gv, 3/0.8/0.2/gvseen, 2/2.2/1/gvseen, 4/2.2/-0.6/gvseen}
            \node[\s] (a\v) at (\x,\y) {\v};
        \draw[gde] (a0) -- (a3); \draw[gde] (a1) -- (a3);
        \draw[gpath, -{Stealth[length=3mm]}] (a3) -- (a2);
        \draw[gpath, -{Stealth[length=3mm]}] (a2) -- (a4);
        \draw[gpath, -{Stealth[length=3mm]}] (a4) -- (a3);
    \end{scope}
    \begin{scope}[xshift=6.8cm]
        \node[font=\small\bfseries] at (2,1.7) {Example 2: $[2,-1,3,1]$};
        \foreach \v/\x in {0/0, 2/1.3, 3/2.6, 1/3.9}
            \node[gv] (b\v) at (\x,0.2) {\v};
        \draw[gde] (b0) -- (b2); \draw[gde] (b2) -- (b3); \draw[gde] (b3) -- (b1);
        \draw[-{Stealth[length=2mm]}, warningred, dashed] (b1) -- (4.8,0.2);
        \node[font=\scriptsize\bfseries, text=warningred, anchor=west] at (4.8,0.2) {$-1$};
    \end{scope}
\end{tikzpicture}
\figcap{Left: a $\rho$-shape; the cycle (orange) has length 3. Right: a single chain that ends in a vertex without an outgoing edge --- no cycle at all.}
\end{diagrambox}

\subsubsection*{2. The graph model}

\begin{modelbox}[Graph Model --- Longest Cycle in a Graph]
\renewcommand{\arraystretch}{1.35}
\begin{tabularx}{\linewidth}{@{}>{\bfseries}l X@{}}
Vertices & $0..n-1$, given directly \\
Edges & $i \to \texttt{edges}[i]$ when $\texttt{edges}[i] \ne -1$; out-degree $\le 1$, stored in the input array (no adjacency list) \\
Weights & none --- a cycle's length is its number of edges = number of vertices \\
Question & over all pieces of this ``functional graph with holes'': the maximum cycle length, $-1$ if no piece has a cycle \\
Size & $|V| \le 10^5$, $|E| \le 10^5$: we need \bigO{n} or \bigO{n \log n}; \bigO{n^2} = $10^{10}$ is far too slow \\
\end{tabularx}
\end{modelbox}

\begin{diagrambox}
\begin{tikzpicture}
    \begin{scope}
        \node[gv] (p0) at (0,0.6) {};
        \node[gv] (p1) at (1.1,0.6) {};
        \node[gvseen] (p2) at (2.2,0.6) {};
        \node[gvseen] (p3) at (3.1,1.2) {};
        \node[gvseen] (p4) at (3.1,0) {};
        \draw[gde] (p0) -- (p1); \draw[gde] (p1) -- (p2);
        \draw[gde] (p2) -- (p3); \draw[gde] (p3) -- (p4); \draw[gde] (p4) -- (p2);
        \node[note] at (1.6,-0.7) {piece with a cycle ($\rho$): \textbf{counts}};
    \end{scope}
    \begin{scope}[xshift=5.6cm]
        \node[gv] (q0) at (0,0.6) {};
        \node[gv] (q1) at (1.1,1.2) {};
        \node[gv] (q2) at (1.1,0) {};
        \node[gvdead] (q3) at (2.2,0.6) {};
        \draw[gde] (q1) -- (q0); \draw[gde] (q2) -- (q0);
        \draw[gde] (q3) -- (q1);
        \draw[-{Stealth[length=2mm]}, warningred, dashed] (q0) -- (-1,0.6);
        \node[note] at (0.6,-0.7) {tree ending in a sink: \textbf{no cycle}};
    \end{scope}
    \node[note, text width=4.2cm, align=left, anchor=west] at (9.3,0.6) {Every piece is one of these two kinds (theory, ``three relatives''). We need the longest cycle over all $\rho$-pieces.};
\end{tikzpicture}
\end{diagrambox}

\subsubsection*{3. Key observations}

\begin{itemize}[leftmargin=*]
    \item \textbf{At most one cycle per piece, and cycles are disjoint.} A vertex has only one outgoing edge, so it can lie on at most one cycle. Hence the total length of all cycles is $\le n$ --- measuring \emph{every} cycle is affordable.
    \item \textbf{A walk finds a cycle iff it returns into itself.} Walking from $s$, we either fall off at $-1$ (no cycle on this walk), reach a vertex seen on an \emph{earlier} walk (its cycle, if any, was measured then), or reach a vertex of the \emph{current} walk (a new cycle).
    \item \textbf{Visit times measure lengths for free.} If vertex $v$ got time $t_v$ when first visited and the walk comes back to $v$ when the clock shows $T$, the cycle has exactly $T - t_v$ vertices.
    \item \textbf{Why not DFS from every vertex?} Each walk can be $\Theta(n)$ long (one long chain into a cycle), so restarting from every vertex costs $\Theta(n^2)$.
\end{itemize}

\subsubsection*{4. Approaches}

\begin{bruteforcebox}
For every vertex $s$, walk at most $n$ steps and check whether we come back to $s$; if we do after \texttt{len} steps, $s$ is on a cycle of length \texttt{len}.

\begin{lstlisting}[style=cpp]
#include <bits/stdc++.h>
using namespace std;

class Solution {
public:
    int longestCycle(vector<int>& edges) {
        int n = edges.size(), best = -1;
        for (int s = 0; s < n; s++) {     // is s on a cycle? walk at most n steps
            int v = edges[s], len = 1;
            while (v != -1 && v != s && len <= n) {
                v = edges[v];
                len++;
            }
            if (v == s) best = max(best, len);   // came back to s after len steps
        }
        return best;
    }
};
\end{lstlisting}

\begin{complexitybox}
\begin{diagrambox}
\begin{tikzpicture}[x=1.15cm]
    \foreach \i/\l in {0/0, 1/1, 2/2, 3/{\dots}, 4/{n{-}3}}
        \node[gv, minimum size=8.5mm, font=\scriptsize\bfseries] (c\i) at (\i,0) {$\l$};
    \node[gvseen, minimum size=8.5mm, font=\scriptsize\bfseries] (c5) at (5.3,0) {$n{-}2$};
    \node[gvseen, minimum size=8.5mm, font=\scriptsize\bfseries] (c6) at (6.6,0) {$n{-}1$};
    \foreach \i [evaluate=\i as \j using int(\i+1)] in {0,...,4} \draw[gde] (c\i) -- (c\j);
    \draw[gde] (c5) to[bend left=35] (c6); \draw[gde] (c6) to[bend left=35] (c5);
    \draw[gback] (c0.south) to[bend right=12] node[lbl, below]{start $0$: about $n$ steps, never returns to $0$} (c6.south);
\end{tikzpicture}
\figcap{The worst case for the brute force: every chain vertex walks to the 2-cycle and around it until the step limit $n$, without ever coming back to itself.}
\end{diagrambox}

\textbf{Time:} \bigO{n^2}: on a chain $0 \to 1 \to \dots \to n{-}3$ ending in a 2-cycle, every start walks $\approx n$ steps before giving up --- $10^{10}$ steps for $n = 10^5$. \textbf{Space:} \bigO{1}. \textbf{Verdict:} \tle{} (used below as the reference on small random graphs).
\end{complexitybox}
\end{bruteforcebox}

\begin{algorithmbox}[Approach 2: One Clock for All Walks]
Keep one global clock and one array \texttt{when[v]} ($0$ = never visited). For every unvisited $s$, remember $\texttt{start} = $ the clock value, then walk, stamping each new vertex with the current time. Stop at $-1$ or at a visited vertex $v$:
\begin{itemize}[leftmargin=*]
    \item $\texttt{when}[v] \ge \texttt{start}$: $v$ was stamped during \emph{this} walk $\Rightarrow$ new cycle of length $\texttt{time} - \texttt{when}[v]$;
    \item $\texttt{when}[v] < \texttt{start}$: an older walk $\Rightarrow$ nothing new.
\end{itemize}
The comparison with \texttt{start} plays exactly the role of ``state 1 versus state 2'' in the theory's Algorithm 2 --- one integer array instead of colours plus a \texttt{path}.

\begin{lstlisting}[style=cpp]
#include <bits/stdc++.h>
using namespace std;

class Solution {
public:
    int longestCycle(vector<int>& edges) {
        int n = edges.size();
        vector<int> when(n, 0);           // 0 = never visited, otherwise the visit time
        int time = 1, best = -1;
        for (int s = 0; s < n; s++) {
            if (when[s] != 0) continue;
            int start = time;             // every vertex of THIS walk gets a time >= start
            int v = s;
            while (v != -1 && when[v] == 0) {
                when[v] = time++;
                v = edges[v];
            }
            if (v != -1 && when[v] >= start)      // closed a cycle inside this walk
                best = max(best, time - when[v]);
        }
        return best;
    }
};
\end{lstlisting}

\begin{complexitybox}
\textbf{Time:} \bigO{n} --- each vertex is stamped once; the \texttt{while} loop runs $n$ times in total over all walks. \textbf{Space:} \bigO{n} for \texttt{when}. \textbf{Verdict:} \ac{}.
\end{complexitybox}
\end{algorithmbox}

\begin{invariantbox}[Why it is correct]
\textbf{Invariant:} during the walk from $s$, the vertices with $\texttt{when} \ge \texttt{start}$ are exactly the vertices visited by this walk, and their times are consecutive: $\texttt{start}, \texttt{start}+1, \dots, \texttt{time}-1$ in walk order.

\emph{New cycle found.} If the walk reaches $v$ with $\texttt{when}[v] \ge \texttt{start}$, then the vertices stamped $\texttt{when}[v], \dots, \texttt{time}-1$ are the walk from $v$ back to (just before) $v$ --- a cycle with $\texttt{time} - \texttt{when}[v]$ vertices.

\emph{No cycle missed.} Every cycle vertex is stamped by some walk. Take the first walk that touches a cycle $C$: it enters $C$ at some vertex and, because no vertex of $C$ was stamped before, it goes all the way round and returns into itself --- so $C$ is measured. Later walks reaching $C$ see old stamps and correctly do nothing. Cycles are disjoint, so each is measured exactly once. \hfill$\blacksquare$
\end{invariantbox}

\begin{alternativebox}[Approach 3: Peel the Trees, Measure What Is Left]
Repeatedly delete vertices of in-degree 0 (Kahn, theory Algorithm 3). Every tree vertex --- including whole sink-trees with no cycle --- is deleted; the survivors are disjoint cycles. Walk once around each survivor cycle, switching its vertices off so that every cycle is measured once.

\begin{lstlisting}[style=cpp]
#include <bits/stdc++.h>
using namespace std;

class Solution {
public:
    int longestCycle(vector<int>& edges) {
        int n = edges.size();
        vector<int> indeg(n, 0);
        for (int v = 0; v < n; v++)
            if (edges[v] != -1) indeg[edges[v]]++;
        queue<int> q;
        for (int v = 0; v < n; v++)
            if (indeg[v] == 0) q.push(v);
        while (!q.empty()) {              // peel vertices nobody points to
            int v = q.front(); q.pop();
            if (edges[v] != -1 && --indeg[edges[v]] == 0) q.push(edges[v]);
        }
        int best = -1;                    // survivors (indeg > 0) form disjoint cycles
        for (int s = 0; s < n; s++) {
            if (indeg[s] == 0) continue;
            int len = 0, v = s;
            do {                          // walk once around, switching vertices off
                indeg[v] = 0;
                v = edges[v];
                len++;
            } while (v != s);
            best = max(best, len);
        }
        return best;
    }
};
\end{lstlisting}

\begin{diagrambox}
\begin{tikzpicture}
    \foreach \v/\x/\y/\s in {0/0/0.8/gvsel, 1/1.3/0/gvsel, 2/0/-0.8/gvsel, 3/3/0.8/gvsel, 4/4.4/0.8/gvsel, 5/4.4/-0.8/gvsel, 6/3/-0.8/gvsel,
                             7/3/2.1/gvdead, 8/6.2/0/gvdead, 9/7.5/0/gvdead}
        \node[\s] (k\v) at (\x,\y) {\v};
    \foreach \u/\w in {0/1, 1/2, 2/0, 3/4, 4/5, 5/6, 6/3} \draw[gde] (k\u) -- (k\w);
    \foreach \u/\w in {7/3, 8/9} \draw[gfaded, dashed, -{Stealth[length=2mm]}] (k\u) -- (k\w);
    \node[gdist, right=1mm of k7] {peeled 1st};
    \node[gdist, above=1mm of k8] {peeled 2nd};
    \node[gdist, below=1mm of k9] {peeled 3rd};
    \node[note, text width=3.9cm, align=left, anchor=west] at (8.3,0.4) {On the trace input (Section 5): peel $7, 8, 9$; the survivors are the cycles $\{0,1,2\}$ (length 3) and $\{3,4,5,6\}$ (length 4) $\Rightarrow$ answer 4.};
\end{tikzpicture}
\end{diagrambox}

\begin{complexitybox}
\textbf{Time:} \bigO{n}. \textbf{Space:} \bigO{n} plus the queue. \textbf{Verdict:} \ac{}. It needs two passes, but it separates the two questions cleanly: ``who is on a cycle?'' and ``how long is each cycle?''.
\end{complexitybox}
\end{alternativebox}

\begin{tricksbox}[Approach 4: Implementation Refinements]
\begin{itemize}[leftmargin=*]
    \item \textbf{Timestamps replace three colours.} The single test $\texttt{when}[v] \ge \texttt{start}$ answers both ``visited?'' and ``visited by this walk?'' --- no \texttt{path} vector, no second colouring pass.
    \item \textbf{Iterative, never recursive.} A recursive DFS along the arrows would go $10^5$ levels deep on a long chain; the \texttt{while} loop has no depth at all.
    \item \textbf{Guard $-1$ first.} Both conditions are written \texttt{v != -1 \&\& \ldots}; with the order reversed, \texttt{when[-1]} is read before the test.
    \item \textbf{No overflow worries:} the clock never exceeds $n + 1$.
\end{itemize}
\end{tricksbox}

\subsubsection*{5. Trace}

Our own input with two cycles of different lengths, a tail into the longer one and a chain without a cycle: \;$\texttt{edges} = [1,2,0,4,5,6,3,3,9,-1]$. Answer: $4$.

\begin{diagrambox}
\begin{tikzpicture}
    \foreach \v/\x/\y in {0/0/0.8, 1/1.3/0, 2/0/-0.8, 3/3/0.8, 4/4.4/0.8, 5/4.4/-0.8, 6/3/-0.8, 7/3/2.1, 8/6.2/0, 9/7.5/0}
        \node[gv] (g\v) at (\x,\y) {\v};
    \foreach \u/\w in {0/1, 1/2, 2/0, 3/4, 4/5, 5/6, 6/3, 7/3, 8/9} \draw[gde] (g\u) -- (g\w);
    \draw[-{Stealth[length=2mm]}, warningred, dashed] (g9) -- (8.5,0);
    \node[font=\scriptsize\bfseries, text=warningred, anchor=west] at (8.5,0) {$-1$};
\end{tikzpicture}
\figcap{The trace graph: cycle $0\to1\to2\to0$, cycle $3\to4\to5\to6\to3$ with tail $7 \to 3$, and the chain $8 \to 9 \to -1$.}
\end{diagrambox}

\begin{tracebox}[{Approach 2 on $[1,2,0,4,5,6,3,3,9,-1]$}]
\centering
\begin{tikzpicture}
    \begin{scope}
        \foreach \v/\x/\y/\s/\t in {0/0/0.8/gvseen/1, 1/1.3/0/gvseen/2, 2/0/-0.8/gvseen/3, 3/3/0.8/gv/{}, 4/4.4/0.8/gv/{}, 5/4.4/-0.8/gv/{}, 6/3/-0.8/gv/{}, 7/3/2.1/gv/{}, 8/6.2/0/gv/{}, 9/7.5/0/gv/{}}{
            \node[\s] (a\v) at (\x,\y) {\v};
            \node[gdist, above left=0mm and 0mm of a\v] {\t};
        }
        \foreach \u/\w in {0/1, 1/2, 3/4, 4/5, 5/6, 6/3, 7/3, 8/9} \draw[gde] (a\u) -- (a\w);
        \draw[gback] (a2) -- (a0);
        \node[note, text width=4.3cm, align=left, anchor=west] at (8.3,0.4) {\textbf{$s=0$}, start $= 1$: stamps $0{:}1$, $1{:}2$, $2{:}3$; clock $= 4$; back at $0$ with $\texttt{when}[0] = 1 \ge 1$ $\Rightarrow$ cycle $4 - 1 = 3$. best $= 3$.};
    \end{scope}
    \begin{scope}[yshift=-3.9cm]
        \foreach \v/\x/\y/\s/\t in {0/0/0.8/gvdone/1, 1/1.3/0/gvdone/2, 2/0/-0.8/gvdone/3, 3/3/0.8/gvseen/4, 4/4.4/0.8/gvseen/5, 5/4.4/-0.8/gvseen/6, 6/3/-0.8/gvseen/7, 7/3/2.1/gv/{}, 8/6.2/0/gv/{}, 9/7.5/0/gv/{}}{
            \node[\s] (b\v) at (\x,\y) {\v};
            \node[gdist, above left=0mm and 0mm of b\v] {\t};
        }
        \foreach \u/\w in {0/1, 1/2, 2/0, 3/4, 4/5, 5/6, 7/3, 8/9} \draw[gde] (b\u) -- (b\w);
        \draw[gback] (b6) -- (b3);
        \node[note, text width=4.3cm, align=left, anchor=west] at (8.3,0.4) {\textbf{$s=3$} ($1, 2$ skipped), start $= 4$: stamps $4..7$; clock $= 8$; back at $3$ with $\texttt{when}[3] = 4 \ge 4$ $\Rightarrow$ cycle $8 - 4 = 4$. best $= 4$.};
    \end{scope}
    \begin{scope}[yshift=-7.8cm]
        \foreach \v/\x/\y/\s/\t in {0/0/0.8/gvdone/1, 1/1.3/0/gvdone/2, 2/0/-0.8/gvdone/3, 3/3/0.8/gvdone/4, 4/4.4/0.8/gvdone/5, 5/4.4/-0.8/gvdone/6, 6/3/-0.8/gvdone/7, 7/3/2.1/gvseen/8, 8/6.2/0/gvseen/9, 9/7.5/0/gvseen/10}{
            \node[\s] (c\v) at (\x,\y) {\v};
            \node[gdist, above left=0mm and 0mm of c\v] {\t};
        }
        \foreach \u/\w in {0/1, 1/2, 2/0, 3/4, 4/5, 5/6, 6/3, 8/9} \draw[gde] (c\u) -- (c\w);
        \draw[gcross] (c7) -- (c3);
        \node[note, text width=4.3cm, align=left, anchor=west] at (8.3,0.4) {\textbf{$s=7$}, start $= 8$: stamp $7{:}8$, reach $3$ with $\texttt{when}[3] = 4 < 8$: old walk, nothing. \textbf{$s=8$}, start $= 9$: stamps $8{:}9$, $9{:}10$, then $-1$: no cycle.};
    \end{scope}
\end{tikzpicture}

\medskip
\small
\begin{tabular}{@{}c|cccccccccc@{}}
\toprule
$v$ & 0 & 1 & 2 & 3 & 4 & 5 & 6 & 7 & 8 & 9 \\
\texttt{edges}$[v]$ & 1 & 2 & 0 & 4 & 5 & 6 & 3 & 3 & 9 & $-1$ \\
\texttt{when}$[v]$ & 1 & 2 & 3 & 4 & 5 & 6 & 7 & 8 & 9 & 10 \\
\midrule
walk ($s$ / start) & 0 / 1 & 0 / 1 & 0 / 1 & 3 / 4 & 3 / 4 & 3 / 4 & 3 / 4 & 7 / 8 & 8 / 9 & 8 / 9 \\
\bottomrule
\end{tabular}
\par\smallskip
Result: cycles of length $3$ and $4$ found, \textbf{answer $4$}.
\end{tracebox}

The two official examples run the same way: in $[3,3,4,2,3]$ the walk from $0$ stamps $0{:}1, 3{:}2, 2{:}3, 4{:}4$ and returns to $3$ (time $2 \ge 1$), giving $5 - 2 = 3$; in $[2,-1,3,1]$ the single walk $0,2,3,1$ falls off at $-1$, so the answer stays $-1$.

\subsubsection*{6. Test suite}

\begin{testbox}
\centering\small
\begin{tabular}{@{}lll@{}}
\toprule
\textbf{edges} & \textbf{Expected} & \textbf{Why this test} \\
\midrule
$[3,3,4,2,3]$ & 3 & official example 1: cycle reached through tails \\
$[2,-1,3,1]$ & $-1$ & official example 2: no cycle at all \\
$[1,0]$ & 2 & smallest cycle (self-loops are forbidden) \\
$[-1,-1]$ & $-1$ & no edges \\
$[1,2,0,4,5,6,3,3,9,-1]$ & 4 & two cycles, tail into the longer one (the trace) \\
$\texttt{edges}[i] = (i+1) \bmod n$, $n = 10^5$ & 100000 & one huge cycle: the maximum answer \\
\bottomrule
\end{tabular}
\end{testbox}

All three programs agree on the examples above and on $200\,000$ random graphs with $2 \le n \le 10$.

\subsubsection*{7. Common mistakes}

\begin{warningbox}[Mistake 1 --- a plain \texttt{visited} array]
\begin{lstlisting}[style=snippet]
if (v != -1 && visited[v]) best = max(best, time - when[v]);   // WRONG
if (v != -1 && when[v] >= start) best = max(best, time - when[v]);   // RIGHT
\end{lstlisting}
With only ``visited / not visited'', the walk from $7$ in the trace hits $3$ (visited earlier) and reports a bogus cycle of length $9 - 4 = 5$. Only a vertex of the \emph{current} walk closes a cycle.
\end{warningbox}

\begin{warningbox}[Mistake 2 --- restarting from every vertex]
Running ``walk and detect'' from all $n$ vertices with a fresh visited set each time is the \bigO{n^2} brute force in disguise. One shared \texttt{when} array makes the whole algorithm \bigO{n}.
\end{warningbox}

\begin{warningbox}[Mistake 3 --- counting the tail into the cycle length]
The cycle length is $\texttt{time} - \texttt{when}[v]$ where $v$ is the \emph{repeated} vertex, not $\texttt{time} - \texttt{start}$: in Example 1 the walk from $0$ has $4$ vertices, but vertex $0$ is a tail vertex and the cycle has only $3$.
\end{warningbox}

\begin{warningbox}[{Mistake 4 --- returning $0$ instead of $-1$}]
Initialise \texttt{best} with $-1$: a graph without cycles must return $-1$, and every real cycle has length $\ge 2$, so $-1$ is never confused with an answer.
\end{warningbox}

\subsubsection*{8. Takeaways}

\begin{itemize}[leftmargin=*]
    \item Out-degree $\le 1$ $\Rightarrow$ cycles are vertex-disjoint and their total length is $\le n$: measure them all.
    \item One global clock: ``visited during this walk'' $\Leftrightarrow$ $\texttt{when}[v] \ge \texttt{start}$; cycle length $= \texttt{time} - \texttt{when}[v]$.
    \item Peeling in-degree-0 vertices removes trees \emph{and} sink-chains, leaving only cycles --- a second, equally linear view.
    \item With arbitrary out-degrees this problem changes character completely (see V5): the ``one arrow'' structure is what makes it easy.
\end{itemize}

\begin{variantsbox}
\variant{Permutation: every vertex on a cycle}{\varin{https://leetcode.com/problems/array-nesting/}}
  {\texttt{nums} is a permutation, so there are no tails and no $-1$.}
  {every walk from an unvisited vertex is exactly one whole cycle.}
  {count the length of each walk with a plain \texttt{seen} array; answer the maximum.}
\variant{First repeated vertex from every start}{\varin{https://codeforces.com/problemset/problem/1020/B}}
  {out-degree exactly 1; for every start output where the walk first repeats.}
  {same walks; the answer is the entry point of each start's $\rho$.}
  {cycle vertices answer themselves, tail vertices copy the answer of their successor (Problem 1.5.1).}
\variant{Number of cycles}{\varnone}
  {count the cycles instead of measuring the longest.}
  {same graph and same walks.}
  {increment a counter whenever $\texttt{when}[v] \ge \texttt{start}$; or count the survivor cycles after peeling.}
\variant{Round table of favourites}{\varin{https://leetcode.com/problems/maximum-employees-to-be-invited-to-a-meeting/}}
  {out-degree exactly 1; seat as many people as possible next to their favourite.}
  {either one long cycle alone, or all 2-cycles together, each extended by the longest chain into both ends.}
  {peeling with ``longest chain ending here'' values; answer $\max(\text{longest cycle},\ \sum_{2\text{-cycles}} \text{chains})$ (E5.7).}
\variant{Any number of outgoing edges}{\varnone}
  {a general directed graph: a vertex may have many outgoing edges.}
  {cycles now overlap and walks are no longer forced.}
  {detecting \emph{a} cycle needs DFS with colours (Chapter~2 (DFS)); the \emph{longest} simple cycle is NP-hard --- no efficient algorithm is known.}
\end{variantsbox}
